% ============================================================
% PROJECT 83: AGENTIC AI COMPILER ASSISTANT
% Comprehensive Project Documentation
% ============================================================
\documentclass[a4paper,11pt]{article}

% --- Packages ---
\usepackage[utf8]{inputenc}
\usepackage[T1]{fontenc}
\usepackage{geometry}
\usepackage{enumitem}
\usepackage{titlesec}
\usepackage[table]{xcolor}
\usepackage{fancyhdr}
\usepackage{tabularx}
\usepackage{listings}
\usepackage{graphicx}
\usepackage{lastpage}
\usepackage{hyperref}
\usepackage{booktabs}
\usepackage{tcolorbox}
\usepackage{amsmath}
\usepackage{array}
\usepackage{mdframed}
\usepackage{tikz}
\usetikzlibrary{shapes.geometric, arrows.meta, positioning}

% --- Page Layout ---
\geometry{top=1in, bottom=1in, left=0.9in, right=0.9in}
\pagestyle{fancy}
\fancyhf{}
\lhead{\small Project 83: Agentic AI Compiler Assistant}
\rhead{\small Comprehensive Project Documentation}
\cfoot{\thepage\ of \pageref{LastPage}}

% --- Section Formatting ---
\titleformat{\section}{\Large\bfseries\color{blue!70!black}}{\thesection.}{1em}{}[\titlerule]
\titleformat{\subsection}{\large\bfseries\color{blue!50!black}}{\thesubsection}{1em}{}
\titleformat{\subsubsection}{\bfseries\color{blue!30!black}}{\thesubsubsection}{1em}{}

% --- Code Listing Style ---
\lstset{
    basicstyle=\ttfamily\footnotesize,
    breaklines=true,
    frame=single,
    backgroundcolor=\color{gray!8},
    keywordstyle=\color{blue!80!black}\bfseries,
    stringstyle=\color{green!50!black},
    commentstyle=\color{gray}\itshape,
    showstringspaces=false,
    tabsize=4,
    captionpos=b,
    numbers=left,
    numberstyle=\tiny\color{gray},
    numbersep=5pt
}

% --- Custom Commands ---
\newcommand{\statusdone}[1]{\colorbox{green!20}{\textbf{\texttt{#1}}}}
\newcommand{\statuspending}[1]{\colorbox{yellow!20}{\textbf{\texttt{#1}}}}
\newcommand{\statusbonus}[1]{\colorbox{purple!15}{\textbf{\texttt{#1}}}}
\newcommand{\gap}[1]{\textit{\textbf{Gap:} #1}}
\newcommand{\solution}[1]{\textit{\textbf{Solution:} #1}}

% ============================================================
\begin{document}
% ============================================================

% -------- TITLE PAGE --------
\begin{titlepage}
\centering
\vspace*{2cm}
{\Huge\bfseries\color{blue!70!black} Project 83}\\[6pt]
{\LARGE\bfseries Conversational Agentic AI Compiler Assistant}\\[10pt]
{\Large Comprehensive Project Documentation}\\[2cm]

\begin{tcolorbox}[colback=blue!5!white, colframe=blue!60!black, width=0.85\textwidth]
\centering
\textbf{A custom-built Mini-Python compiler front-end, integrated with\\
Google Gemini 2.5 Flash AI, that provides conversational,\\
context-aware error explanations and code fix suggestions.}
\end{tcolorbox}

\vspace{1.5cm}
\begin{tabular}{ll}
\toprule
\textbf{Project ID}     & Project 83 \\
\textbf{Target Language}& Mini-Python (Python Subset) \\
\textbf{AI Engine}      & Google Gemini 2.5 Flash (via \texttt{google-generativeai}) \\
\textbf{Language}       & Python 3.12 \\
\textbf{Current Phase}  & Week 8 --- Gemini Integration \& Context Injection \\
\textbf{Test Count}     & 128 tests --- All Passing \\
\textbf{Status Date}    & February 24, 2026 \\
\bottomrule
\end{tabular}

\vfill
{\large\textit{Pipeline:} Source Code $\rightarrow$ Lexer $\rightarrow$ Parser $\rightarrow$ Semantic Analyzer $\rightarrow$ Context Provider $\rightarrow$ Gemini AI}
\end{titlepage}

% -------- TABLE OF CONTENTS --------
\tableofcontents
\newpage

% ============================================================
\section{Project Overview and Motivation}
% ============================================================

\subsection{Problem Statement}

Traditional compilers, even modern ones like CPython, produce error messages that are:
\begin{itemize}
    \item \textbf{Cryptic} for beginners and intermediate programmers
    \item \textbf{Context-free} --- they report \textit{what} went wrong but not \textit{why} or \textit{how} to fix it
    \item \textbf{Security-blind} --- they cannot identify code that is syntactically valid but logically dangerous (e.g., using \texttt{eval()}, command injection patterns)
    \item \textbf{Non-conversational} --- the developer gets a one-time error message with no ability to ask follow-up questions
\end{itemize}

\subsection{Project Vision}

Project 83 addresses these gaps by building a \textbf{Conversational Agentic AI Compiler Assistant}. The project transforms a standard compiler workflow --- where the tool simply reports errors and stops --- into an \textbf{interactive debugging partnership} between the developer and an AI agent that has deep knowledge of the code's structure.

The system is built in two parallel tracks:
\begin{enumerate}
    \item \textbf{Compiler Front-End}: A complete, hand-crafted compiler pipeline for a Python subset (Mini-Python) built from scratch --- no compiler-generators (ANTLR, PLY, etc.) are used.
    \item \textbf{Agentic AI Layer}: Google Gemini 2.5 Flash is wrapped with custom orchestration logic that feeds it the compiler's internal state (AST, Symbol Table, Errors) as structured context, enabling grammar-grounded, intelligent explanations.
\end{enumerate}

\subsection{What Was Needed vs. What Was Delivered}

\begin{table}[h!]
\centering
\begin{tabularx}{\textwidth}{|l|X|X|}
\hline
\rowcolor{blue!10}
\textbf{Week} & \textbf{What Was Needed} & \textbf{What Was Delivered} \\
\hline
1 & Project definition, EBNF grammar design, agent interaction flow planning & Grammar file (\texttt{mini\_python.ebnf}), project structure, interaction flow diagrams \\
\hline
2 & Gemini API setup, prompt architecture, basic AI test cases & Working \texttt{gemini\_client.py}, \texttt{.env} key management, AI test cases \\
\hline
3 & SRS document, security threat model & Software Requirements Spec (SRS), threat model \\
\hline
4 & Architecture design, module interaction diagrams, tech stack justification & Architecture blueprint, module interaction docs, compiler directory structure \\
\hline
5 & Regex-based Lexer for Mini-Python & \texttt{lexer.py} with indentation handling, 13 unit tests passing \\
\hline
6 & Recursive Descent Parser, AST construction & \texttt{parser.py} + \texttt{ast\_nodes.py} with 16 node types, 27 tests passing \\
\hline
7 & Symbol Table, scope analysis, type checking, diagnostics & \texttt{semantic\_analyzer.py} with full scope stack, 16 tests passing \\
\hline
8 & Gemini Context Provider, secure API integration & \texttt{context\_provider.py} + enhanced \texttt{gemini\_client.py}, 27 deterministic tests. \textbf{Bonus:} TAC generator (39 tests) \\
\hline
\end{tabularx}
\caption{Planned objectives vs. delivered results by week}
\end{table}


% ============================================================
\section{Technologies and Tools Used}
% ============================================================

\subsection{Core Programming Language and Runtime}

\begin{itemize}
    \item \textbf{Python 3.12}: The primary implementation language. Used for all compiler components and the AI orchestration layer.
    \item \textbf{Python \texttt{re} module}: Powers the regex-based lexer tokenizer.
    \item \textbf{Python \texttt{enum} module}: Defines the \texttt{TokenType} and \texttt{SymbolType} enumerations for type safety.
    \item \textbf{Python \texttt{dataclasses}}: Used to define all 16 AST node classes cleanly and concisely.
    \item \textbf{Python \texttt{typing}}: Provides type hints throughout all modules for maintainability.
    \item \textbf{Python \texttt{logging}}: Used in the Gemini client for structured logging (replacing debug print statements).
\end{itemize}

\subsection{AI / Machine Learning Stack}

\begin{itemize}
    \item \textbf{Google Gemini 2.5 Flash (\texttt{models/gemini-2.5-flash})}: The large language model (LLM) used as the agent's reasoning engine. Accessed via the official Python SDK.
    \item \textbf{\texttt{google-generativeai}}: Official Google Python client library for the Gemini API. Provides model instantiation, chat sessions, function calling, and safety filter management.
    \item \textbf{System Prompt Injection}: The Gemini model is initialized with a persistent system instruction containing the project's EBNF grammar, agent role definition, and ReAct reasoning rules.
    \item \textbf{Context Injection (RAG-like)}: The compiler's internal state is serialized into natural language and injected into every Gemini prompt, acting as a retrieval-augmented context layer.
    \item \textbf{ReAct (Reason + Act) Pattern}: The agent is instructed to first reason about the error against the grammar before producing a response or suggesting a fix.
\end{itemize}

\subsection{Development and Testing Tools}

\begin{itemize}
    \item \textbf{pytest}: Testing framework for all 128 unit and integration tests.
    \item \textbf{python-dotenv}: Secure API key management via \texttt{.env} file.
    \item \textbf{LaTeX (Overleaf)}: All project reports and documentation written in professional LaTeX.
    \item \textbf{Draw.io / LucidChart}: Used for architecture and interaction flow diagrams during planning phases.
\end{itemize}

\subsection{Grammar Specification}

The project uses \textbf{EBNF (Extended Backus-Naur Form)} to formally define the Mini-Python grammar targeted by the compiler. This grammar file (\texttt{grammar/mini\_python.ebnf}) serves as the single source-of-truth for both the parser implementation and the AI agent's reasoning.

\begin{lstlisting}[language=, caption=Mini-Python EBNF Grammar (grammar/mini\_python.ebnf)]
program         = { statement } ;
statement       = assignment | if_stmt | while_stmt | func_def | print_stmt ;
assignment      = identifier , "=" , expression , NEWLINE ;
if_stmt         = "if" , condition , ":" , block , [ "else" , ":" , block ] ;
while_stmt      = "while" , condition , ":" , block ;
func_def        = "def" , identifier , "(" , [ params ] , ")" , ":" , block ;
print_stmt      = "print" , "(" , expression , ")" , NEWLINE ;
block           = NEWLINE , INDENT , { statement } , DEDENT ;
expression      = term , { ( "+" | "-" ) , term } ;
term            = factor , { ( "*" | "/" ) , factor } ;
factor          = identifier | number | "(" , expression , ")" ;
condition       = expression , comp_op , expression ;
comp_op         = "==" | "!=" | "<" | ">" | "<=" | ">=" ;
identifier      = letter , { letter | digit | "_" } ;
number          = digit , { digit } ;
\end{lstlisting}


% ============================================================
\section{System Architecture}
% ============================================================

\subsection{Three-Layer Architecture}

The system is organized into three distinct architectural layers, each with a clear responsibility:

\begin{tcolorbox}[colback=blue!3!white, colframe=blue!60!black, title={System Layers Overview}]
\begin{description}[style=nextline]
    \item[Layer 1 --- Deterministic Compiler Front-End]
    Hand-crafted lexer, parser, and semantic analyzer. Produces the AST, Symbol Table, and error diagnostics. No AI involved in this layer.
    
    \item[Layer 2 --- Agentic Orchestration Layer (Gemini AI)]
    The Context Provider serializes compiler output into natural language. The Gemini client sends this context to the AI with a system prompt embedding the grammar. The AI reasons using ReAct patterns.
    
    \item[Layer 3 --- Conversational Interface (Planned)]
    A Streamlit web UI or Rich CLI allowing users to submit code, see diagnostics, and engage in a follow-up conversation with the AI agent.
\end{description}
\end{tcolorbox}

\subsection{Data Flow Pipeline}

The end-to-end compilation and AI analysis pipeline can be described as:

\[
\text{Source Code} \xrightarrow{\text{Lexer}} \text{Tokens} \xrightarrow{\text{Parser}} \text{AST} \xrightarrow{\text{Semantic Analyzer}} \text{Symbols + Errors} \xrightarrow{\text{Context Provider}} \text{Natural Language} \xrightarrow{\text{Gemini}} \text{Explanation}
\]

\begin{table}[h!]
\centering
\begin{tabularx}{\textwidth}{|c|l|l|X|}
\hline
\rowcolor{blue!10}
\textbf{Phase} & \textbf{Component} & \textbf{File} & \textbf{Output} \\
\hline
1 & Tokenizer & \texttt{lexer.py} & Stream of typed Tokens \\
\hline
2 & AST Builder & \texttt{parser.py} & Hierarchical AST (Program node) \\
\hline
3 & Symbol Table & \texttt{semantic\_analyzer.py} & Scope-annotated symbol records \\
\hline
4 & Diagnostics & \texttt{semantic\_analyzer.py} & SemanticError / Warning list \\
\hline
5 & Context Serialization & \texttt{context\_provider.py} & Natural language context string \\
\hline
6 & AI Reasoning & \texttt{gemini\_client.py} & Conversational error explanation \\
\hline
\end{tabularx}
\caption{Complete pipeline phase-by-phase}
\end{table}

\subsection{Directory Structure}

\begin{lstlisting}[language=, caption=Project Directory Layout]
CD-Project/
|--- src/
|    |--- compiler/
|    |    |--- tokens.py           # TokenType enum + Token class
|    |    |--- lexer.py            # Regex-based Lexer + indentation tracking
|    |    |--- ast_nodes.py        # 16 AST node dataclasses
|    |    |--- parser.py           # Recursive Descent Parser
|    |    |--- semantic_analyzer.py # Symbol Table + Scope + Type checks
|    |    `--- tac_generator.py    # [BONUS] Three-Address Code generator
|    `--- orchestrator/
|         |--- context_provider.py # Compiler state -> natural language
|         `--- gemini_client.py    # Gemini API wrapper + system prompt
|--- tests/
|    |--- test_week5_lexer.py      # 13 Lexer tests
|    |--- test_week6_parser.py     # 27 Parser tests
|    |--- test_week7_semantic.py   # 16 Semantic tests
|    |--- test_week8_context.py    # 27 Context Provider tests
|    |--- test_week8_tac.py        # 39 TAC Generator tests [BONUS]
|    `--- test_integration_weeks5_7.py  # 6 integration tests
|--- grammar/
|    `--- mini_python.ebnf         # Formal grammar specification
|--- docs/
|    |--- ProjectDescriptionMethodologicalApproach.tex
|    |--- DetailedExecutionPlanandReport.tex
|    `--- prompts.md
|--- reports/
|    |--- week_01_report.tex ... week_08_report.tex
|    |--- SRS_latex.tex
|    `--- comprehensive_project_documentation.tex
|--- demo_week5_lexer.py   ... demo_week8_tac.py
|--- requirements.txt
`--- .env                          # GEMINI_API_KEY (not committed)
\end{lstlisting}


% ============================================================
\section{Implementation Deep-Dive}
% ============================================================

\subsection{Week 5: Lexical Analyzer (Lexer)}

\subsubsection{What It Does}

The Lexer (\texttt{src/compiler/lexer.py}) is the first stage of compilation. It converts raw Mini-Python source code into a flat sequence of typed \textbf{Token} objects.

\subsubsection{Key Design Decisions}

\begin{itemize}
    \item \textbf{Regex-Based Scanning}: A single compiled regex pattern (using alternation with named groups \texttt{(?P<NAME>pattern)}) matches every token type in a single pass using Python's \texttt{re.match}.
    \item \textbf{Python-Specific Indentation}: Python's significant whitespace (indentation as block delimiters) is handled explicitly. An \texttt{indent\_stack} tracks indentation levels across lines, emitting \texttt{INDENT} and \texttt{DEDENT} tokens to mark block boundaries.
    \item \textbf{Keyword Detection}: Identifiers are first regex-matched, then checked against a keyword map (e.g., \texttt{if}, \texttt{def}, \texttt{while}, \texttt{return}, \texttt{print}) to produce the correct keyword token type.
    \item \textbf{Line and Column Tracking}: Every token stores its \texttt{line} and \texttt{column} numbers for downstream error reporting.
\end{itemize}

\subsubsection{Token Types Implemented}

\begin{table}[h!]
\centering
\begin{tabular}{|l|l|}
\hline
\rowcolor{blue!10}
\textbf{Category} & \textbf{Tokens} \\
\hline
Special & \texttt{EOF, NEWLINE, INDENT, DEDENT} \\
\hline
Literals & \texttt{NUMBER, STRING, IDENTIFIER} \\
\hline
Keywords & \texttt{DEF, IF, ELSE, WHILE, RETURN, PRINT} \\
\hline
Arithmetic & \texttt{PLUS, MINUS, MULTIPLY, DIVIDE} \\
\hline
Comparison & \texttt{EQ, NEQ, LT, GT, LTE, GTE} \\
\hline
Assignment & \texttt{ASSIGN} \\
\hline
Delimiters & \texttt{LPAREN, RPAREN, COLON, COMMA} \\
\hline
\end{tabular}
\caption{Token types defined in \texttt{tokens.py}}
\end{table}

\subsection{Week 6: Syntax Parser (Recursive Descent)}

\subsubsection{What It Does}

The Parser (\texttt{src/compiler/parser.py}) takes the token stream from the Lexer and builds a hierarchical \textbf{Abstract Syntax Tree (AST)}, validating that the code conforms to the EBNF grammar.

\subsubsection{Key Design Decisions}

\begin{itemize}
    \item \textbf{Recursive Descent Parsing}: A hand-written top-down parser is used, with one parsing function per grammar rule. This approach was chosen over parser generators (like ANTLR or PLY) to provide maximum control and academic rigor.
    \item \textbf{Operator Precedence}: Arithmetic precedence (\texttt{*} before \texttt{+}) is naturally encoded in the grammar's production hierarchy: \texttt{expression} $\rightarrow$ \texttt{term} $\rightarrow$ \texttt{factor}.
    \item \textbf{Unary Minus Support}: The \texttt{parse\_factor()} function was extended to handle the \texttt{-expr} pattern, producing a \texttt{UnaryOp} node.
    \item \textbf{Visitor Pattern}: All AST nodes implement an \texttt{accept(visitor)} method to support future traversals (used in the semantic analyzer and context provider).
\end{itemize}

\subsubsection{AST Node Hierarchy (16 Node Types)}

\begin{table}[h!]
\centering
\begin{tabular}{|l|l|l|}
\hline
\rowcolor{blue!10}
\textbf{Category} & \textbf{Node Class} & \textbf{Description} \\
\hline
Root & \texttt{Program} & Root node, contains all top-level statements \\
\hline
Container & \texttt{Block} & Indented sequence of statements \\
\hline
\multirow{7}{*}{Statements} & \texttt{Assignment} & \texttt{x = expr} \\
 & \texttt{IfStatement} & \texttt{if cond: ... else: ...} \\
 & \texttt{WhileStatement} & \texttt{while cond: ...} \\
 & \texttt{FunctionDef} & \texttt{def name(params): ...} \\
 & \texttt{ReturnStatement} & \texttt{return expr} \\
 & \texttt{PrintStatement} & \texttt{print(expr)} \\
 & \texttt{Statement} & Abstract base class \\
\hline
\multirow{6}{*}{Expressions} & \texttt{BinaryOp} & \texttt{a + b}, \texttt{x == y} \\
 & \texttt{UnaryOp} & \texttt{-x} \\
 & \texttt{Identifier} & Variable/function name reference \\
 & \texttt{Literal} & Integer or string literal \\
 & \texttt{FunctionCall} & \texttt{func(arg1, arg2)} \\
 & \texttt{Expression} & Abstract base class \\
\hline
\multirow{1}{*}{Base} & \texttt{ASTNode} & Base for all nodes (line, column) \\
\hline
\end{tabular}
\caption{Complete AST node hierarchy (16 types in \texttt{ast\_nodes.py})}
\end{table}

\subsection{Week 7: Semantic Analyzer}

\subsubsection{What It Does}

The Semantic Analyzer (\texttt{src/compiler/semantic\_analyzer.py}) traverses the AST and performs meaning-level validation that the parser cannot catch.

\subsubsection{Key Design Decisions}

\begin{itemize}
    \item \textbf{Dictionary-Based Symbol Table}: Each \texttt{Scope} object is a dictionary mapping variable names to \texttt{Symbol} records. This provides O(1) average-case lookup and is straightforward to serialize.
    \item \textbf{Scope Stack}: A \texttt{SymbolTable} manages a list of \texttt{Scope} objects. Function calls push a new scope; returning pops it. Lookup searches from innermost to outermost scope (shadowing semantics).
    \item \textbf{Graceful Error Accumulation}: Errors are accumulated in a list rather than raising on first error, allowing the analyzer to report all errors in a single pass.
\end{itemize}

\subsubsection{Semantic Checks Implemented}

\begin{table}[h!]
\centering
\begin{tabularx}{\textwidth}{|l|X|}
\hline
\rowcolor{blue!10}
\textbf{Check} & \textbf{Description} \\
\hline
Undefined Variable & Error when a variable is used before assignment \\
\hline
Undefined Function & Error when a function is called before its \texttt{def} \\
\hline
Duplicate Definition & Error when a function/variable is redefined in the same scope \\
\hline
Duplicate Parameters & Error when a function parameter name appears twice \\
\hline
Type Mismatch Warning & Warning when arithmetic operands have incompatible types \\
\hline
Type Reassignment Warning & Warning when a variable's type changes on reassignment \\
\hline
Non-callable Call & Error when a non-function identifier is called as a function \\
\hline
Unary Minus Type & Warning when unary minus is applied to a non-integer \\
\hline
\end{tabularx}
\caption{Semantic validation checks implemented}
\end{table}

\subsection{Week 8: Gemini API Integration \& Context Injection}

\subsubsection{Context Provider (\texttt{context\_provider.py})}

The \texttt{ContextProvider} class is the bridge between the compiler and the AI. It walks the compiler's data structures and converts them to structured natural language.

\begin{lstlisting}[language=Python, caption=ContextProvider API]
class ContextProvider:
    @staticmethod
    def serialize_ast(ast: Program) -> str:
        """Convert the AST into a human-readable program structure summary."""

    @staticmethod
    def serialize_symbol_table(symbol_table: SymbolTable) -> str:
        """Convert the scope stack + symbols into readable text."""

    @staticmethod
    def serialize_errors(errors, warnings) -> str:
        """Format semantic errors/warnings for AI consumption."""

    @staticmethod
    def build_context(source_code, ast, analyzer) -> str:
        """Main entry point: combine all above into one context string."""

    @staticmethod
    def build_error_context(source_code, error_type, error_message, ...) -> str:
        """For lexer/parser errors that occur before the AST exists."""
\end{lstlisting}

\subsubsection{Enhanced Gemini Client (\texttt{gemini\_client.py})}

\begin{lstlisting}[language=Python, caption=GeminiCompilerAgent Key Methods]
class GeminiCompilerAgent:
    def analyze_with_context(self, compiler_context: str) -> str:
        """Send full compiler state (AST + Symbols + Errors) to Gemini."""

    def explain_error(self, source_code, error_type, error_message,
                      compiler_context=None) -> str:
        """Ask Gemini to explain a specific compilation error."""

    def _safe_send(self, prompt: str) -> str:
        """Internal: handles rate limiting, timeouts, safety filters."""
\end{lstlisting}

\subsubsection{Bonus: Three-Address Code (TAC) Generator}

As a bonus deliverable, \texttt{tac\_generator.py} implements an intermediate code generator that emits Three-Address Code (TAC) instructions from the AST. This is a standard compiler middle-end phase (between front-end analysis and back-end code generation). The TAC generator supports 21 operation codes including arithmetic operations, conditional jumps, function calls, and return statements.


% ============================================================
\section{Research and Academic References}
% ============================================================

\subsection{Primary Compiler Theory References}

The following foundational textbooks and resources directly informed the implementation choices in this project:

\begin{enumerate}
    \item \textbf{Aho, Lam, Sethi, and Ullman --- ``Compilers: Principles, Techniques, and Tools'' (Dragon Book, 2nd Ed.)}
    \begin{itemize}
        \item \textit{Applied to}: Lexer design (Chapter 3 --- Lexical Analysis), regex-based tokenizers, symbol table construction, scope management
        \item \textit{Key concepts used}: DFA-based scanning, symbol table organization, scope rules
    \end{itemize}

    \item \textbf{Robert Nystrom --- ``Crafting Interpreters''}
    \begin{itemize}
        \item \textit{Applied to}: Recursive descent parsing (Weeks 6--7), AST node design, visitor pattern implementation
        \item \textit{Key concepts used}: Top-down parsing, operator precedence via grammar hierarchy, tree-walking interpreters
    \end{itemize}

    \item \textbf{Andrew Appel --- ``Modern Compiler Implementation in ML''}
    \begin{itemize}
        \item \textit{Applied to}: Symbol table design, semantic analysis rules, scope stack implementation (Week 7)
        \item \textit{Key concepts used}: Environment-based scope stacks, type systems for dynamic languages
    \end{itemize}
\end{enumerate}

\subsection{AI and Agentic System References}

\begin{enumerate}
    \item \textbf{Yao et al. --- ``ReAct: Synergizing Reasoning and Acting in Language Models'' (2022)}
    \begin{itemize}
        \item \textit{Applied to}: The Gemini system prompt instructs the agent to follow the ReAct (Reason + Act) pattern: first reason about the compiler context, then act (suggest a fix).
        \item \textit{Key concepts used}: Chain-of-thought reasoning, interleaving reasoning and action steps
    \end{itemize}

    \item \textbf{Google --- Gemini API Documentation and Function Calling Guide}
    \begin{itemize}
        \item \textit{Applied to}: Gemini model configuration, system prompt injection, safety filter handling, model listing and fallback logic
        \item \textit{Key concepts used}: \texttt{GenerativeModel} instantiation, \texttt{start\_chat()} session management, function calling definitions (planned for Week 9)
    \end{itemize}

    \item \textbf{DeepLearning.AI --- ``Prompt Engineering for Developers''}
    \begin{itemize}
        \item \textit{Applied to}: System prompt design (Week 8), instructing the model to ground responses in the EBNF grammar
        \item \textit{Key concepts used}: Role assignment, chain-of-thought elicitation, structured output prompting
    \end{itemize}

    \item \textbf{NVIDIA Technical Blog --- ``Building LLM-Powered Agents''}
    \begin{itemize}
        \item \textit{Applied to}: Architecture design of the Compiler-Agent Bridge (Week 4 planning), JSON/text serialization of internal state for LLM consumption
        \item \textit{Key concepts used}: Tool-use agents, context window management, stateful agent sessions
    \end{itemize}
\end{enumerate}

\subsection{Security References}

\begin{enumerate}
    \item \textbf{OWASP Top 10 for LLM Applications (2023)}
    \begin{itemize}
        \item \textit{Applied to}: Week 3 --- Software Requirements Specification security section; Week 11 (planned) --- Security Auditing Agent
        \item \textit{Key risks addressed}: Prompt injection, insecure output handling, training data poisoning
    \end{itemize}

    \item \textbf{Python Security Best Practices}
    \begin{itemize}
        \item \textit{Applied to}: Detection of dangerous patterns (\texttt{eval}, \texttt{exec}, \texttt{\_\_import\_\_}) flagged by the Gemini system prompt
        \item \textit{Key concepts used}: Static analysis security testing (SAST) principles
    \end{itemize}
\end{enumerate}


% ============================================================
\section{Gaps Identified and Solutions Implemented}
% ============================================================

\subsection{Gap 1: Static Error Messages with No Contextual Intelligence}

\gap{Traditional compilers report errors as static one-liners (e.g., \texttt{SyntaxError: invalid syntax}). These messages lack root-cause analysis, grammar rule references, and suggested fixes.}

\solution{The \texttt{ContextProvider} serializes the AST, Symbol Table, and diagnostics into a structured natural language context. This context is injected into the Gemini prompt, enabling the AI to provide grammar-grounded explanations that reference specific line numbers and EBNF rules.}

\subsection{Gap 2: LLMs Cannot Reason About Internal Compiler State}

\gap{Sending raw source code to an LLM gives it no access to the compiler's internal understanding of the program --- the variable types, scope boundaries, or the parsed tree structure.}

\solution{The Context Provider converts the compiler's \textit{internal data structures} (AST nodes, Scope objects, Symbol records) into readable text. When Gemini receives ``\texttt{SYMBOL TABLE: Scope: global (level 1), factorial: type=function, defined at line 1}'', it can reason much more precisely than from source code alone.}

\subsection{Gap 3: Non-Deterministic Testing of AI-Integrated Systems}

\gap{Tests that make live API calls to Gemini are slow, expensive, and non-deterministic (Gemini's outputs vary run-to-run). This makes CI/CD integration and repeatable test suites impossible.}

\solution{The Context Provider layer is fully decoupled from the Gemini API. All 27 Week 8 tests in \texttt{test\_week8\_context.py} test only the serialization logic (verifying the context string contains expected keywords, structure, and line numbers), with zero API calls. The optional API call is gated behind a \texttt{--api} flag in the demo scripts.}

\subsection{Gap 4: Python's Significant Whitespace in Hand-Crafted Lexers}

\gap{Python's indentation-based block structure is notoriously difficult to handle in a scanner because indentation is semantically significant --- a change in indent level marks the start or end of a block. Standard character-by-character scanners cannot handle this without explicit state management.}

\solution{An \texttt{indent\_stack} (a list of integer indent levels) is maintained across lines. When a new line's indent level exceeds the top of the stack, an \texttt{INDENT} token is emitted; when it decreases, one or more \texttt{DEDENT} tokens are emitted to close all affected blocks.}

\subsection{Gap 5: Backward Compatibility During Iterative Enhancement}

\gap{The Gemini client was first implemented in Week 2 as a simple wrapper. Week 8 required significant enhancement (system prompts, context-aware methods, secure error handling), risking breaking existing demos and tests.}

\solution{New methods (\texttt{analyze\_with\_context()}, \texttt{explain\_error()}) were added alongside the existing \texttt{get\_assistance()} method, which was kept unchanged. The internal \texttt{\_safe\_send()} helper centralizes all error handling, used by all methods without duplicating logic.}

\subsection{Gap 6: AST-to-Text Fidelity Loss}

\gap{Converting a recursive tree structure to text risks losing structural information (nesting, node types, line positions) that the AI needs to reason accurately.}

\solution{The AST serializer uses hierarchical indentation with per-node-type description methods. The output explicitly includes node type labels, parameter names, line numbers, and nesting level, ensuring the AI receives a faithful structural overview of the program.}


% ============================================================
\section{Testing Strategy and Results}
% ============================================================

\subsection{Test Suite Structure}

\begin{table}[h!]
\centering
\begin{tabularx}{\textwidth}{|l|c|c|X|}
\hline
\rowcolor{blue!10}
\textbf{Test File} & \textbf{Tests} & \textbf{Week} & \textbf{Coverage} \\
\hline
\texttt{test\_week5\_lexer.py} & 13 & 5 & Token types, keywords, operators, indentation, strings, errors \\
\hline
\texttt{test\_week6\_parser.py} & 27 & 6 & Assignments, if/else, while, function def, return, print, expressions \\
\hline
\texttt{test\_week7\_semantic.py} & 16 & 7 & Undefined vars, duplicate defs, type checks, scope management \\
\hline
\texttt{test\_week8\_context.py} & 27 & 8 & AST serialization, symbol table text, error formatting, full pipeline \\
\hline
\texttt{test\_week8\_tac.py} & 39 & 8 (Bonus) & TAC generation for all statement types, 21 operation codes \\
\hline
\texttt{test\_integration\_weeks5\_7.py} & 6 & 5--7 & Full pipeline (Lexer $\rightarrow$ Parser $\rightarrow$ Semantic) integration \\
\hline
\textbf{Grand Total} & \textbf{128} & \textbf{5--8} & \textbf{All Passing} \\
\hline
\end{tabularx}
\caption{Complete test suite --- 128 tests, all passing as of Week 8}
\end{table}

\subsection{Testing Principles}

\begin{itemize}
    \item \textbf{Deterministic Only}: AI API calls are excluded from all automated tests. The Context Provider test suite verifies serialization outputs without requiring Gemini availability.
    \item \textbf{Isolated Unit Tests}: Each compiler stage is tested independently before integration tests verify the end-to-end pipeline.
    \item \textbf{Error Case Coverage}: Tests explicitly verify that the system correctly detects and reports errors (not just valid code).
    \item \textbf{Run Command}: \texttt{python -m pytest tests/ -v} (run from project root)
\end{itemize}

\subsection{Demonstration Scripts}

\begin{table}[h!]
\centering
\begin{tabular}{|l|l|}
\hline
\rowcolor{blue!10}
\textbf{Script} & \textbf{What It Demonstrates} \\
\hline
\texttt{demo\_week5\_lexer.py} & Tokenization of sample Mini-Python code \\
\hline
\texttt{demo\_week6\_parser.py} & AST construction and pretty-print output \\
\hline
\texttt{demo\_week7\_semantic.py} & Symbol table population and error detection \\
\hline
\texttt{demo\_week8\_ai.py} & Context Provider + (optional) live Gemini call (\texttt{--api} flag) \\
\hline
\texttt{demo\_week8\_tac.py} & Three-Address Code generation from source \\
\hline
\end{tabular}
\caption{Demo scripts for each project phase}
\end{table}


% ============================================================
\section{Security Considerations}
% ============================================================

\subsection{API Key Security}

\begin{itemize}
    \item The \texttt{GEMINI\_API\_KEY} is stored exclusively in the \texttt{.env} file (excluded from version control via \texttt{.gitignore}).
    \item The \texttt{python-dotenv} library loads the key at runtime; it is \textbf{never} hardcoded in source files.
    \item The Gemini client trims the key and logs only the first 5 and last 4 characters for verification without exposing the full key in logs.
\end{itemize}

\subsection{AI Safety Filter Handling}

The \texttt{\_safe\_send()} method explicitly catches \texttt{BlockedPromptException}, returning a user-friendly message when safety filters are triggered rather than crashing the application.

\subsection{Security Auditing (In System Prompt)}

The Gemini system prompt instructs the AI agent to flag code patterns that are syntactically valid but security-dangerous:
\begin{itemize}
    \item Use of \texttt{eval()}, \texttt{exec()}, or \texttt{\_\_import\_\_()}
    \item Untrusted input being passed into sensitive sinks
    \item Insecure default patterns
\end{itemize}

\subsection{Planned: Security Agent (Week 11)}

A dedicated Security Auditing Agent will perform deeper static analysis (SAST) on the AST, checking for ``Logic Bombs'', ``Prompt Injection'' within code comments, and enforcing security policies at compile time. This is aligned with the OWASP ``Shift-Left'' security methodology.


% ============================================================
\section{Project Timeline and Status}
% ============================================================

\subsection{Completed Phases (Weeks 1--8)}

\begin{table}[h!]
\centering
\begin{tabularx}{\textwidth}{|c|l|X|c|}
\hline
\rowcolor{blue!10}
\textbf{Week} & \textbf{Phase} & \textbf{Key Deliverable} & \textbf{Status} \\
\hline
1 & Problem Definition & EBNF grammar, project structure & \statusdone{DONE} \\
\hline
2 & API Design & Gemini client, prompt architecture & \statusdone{DONE} \\
\hline
3 & Requirements & SRS document, threat model & \statusdone{DONE} \\
\hline
4 & Architecture & Component blueprint, module interaction maps & \statusdone{DONE} \\
\hline
5 & Lexer & Regex tokenizer, indentation handling (13 tests) & \statusdone{DONE} \\
\hline
6 & Parser & Recursive descent, AST (27 tests) & \statusdone{DONE} \\
\hline
7 & Semantic Analysis & Symbol table, type checking (16 tests) & \statusdone{DONE} \\
\hline
8 & AI Bridge & Context Provider + enhanced Gemini client (27 tests + 39 bonus TAC tests) & \statusdone{DONE} \\
\hline
\end{tabularx}
\caption{Weeks 1--8 status (all complete)}
\end{table}

\subsection{Planned Phases (Weeks 9--16)}

\begin{table}[h!]
\centering
\begin{tabularx}{\textwidth}{|c|l|X|}
\hline
\rowcolor{yellow!15}
\textbf{Week} & \textbf{Phase} & \textbf{Planned Deliverable} \\
\hline
9 & Agentic Tool-Use & Gemini Function Calling --- \texttt{check\_scope()}, \texttt{reparse\_code()} \\
\hline
10 & Conversational UI & Chatbot CLI/Web with follow-up question support \\
\hline
11 & Security Agent & AST-level security auditing, logic bomb detection \\
\hline
12 & Benchmarking & Latency measurement, accuracy vs standard error messages \\
\hline
13 & Explainability (XAI) & Reasoning log showing which AST node led to AI decision \\
\hline
14 & Integration Testing & End-to-end: AI-fixed code re-parsed through the compiler \\
\hline
15 & Documentation & Final report, architecture diagrams, prompt templates \\
\hline
16 & Demo \& Submission & Live demo recording, GitHub submission \\
\hline
\end{tabularx}
\caption{Weeks 9--16 planned phases}
\end{table}


% ============================================================
\section{Key Achievements and Innovation}
% ============================================================

\subsection{Technical Innovations}

\begin{enumerate}
    \item \textbf{Grammar-Grounded AI Reasoning}: By injecting the EBNF grammar into the system prompt, Gemini can reference specific grammar rules when explaining errors --- a level of precision not achievable with generic LLM prompts.

    \item \textbf{Compiler-State Context Injection}: The Context Provider pattern (serializing AST + Symbol Table + Errors as natural language) is a novel bridge between formal compiler theory and modern LLM prompting, producing a form of ``Retrieval-Augmented Generation'' from the compiler's own output.

    \item \textbf{From-Scratch Compiler}: Building the lexer, parser, and semantic analyzer without compiler-generators ensures the system has deep, controllable access to its internal state --- a requirement for meaningful AI integration.

    \item \textbf{Deterministic AI Testing}: Decoupling the Context Provider from the API layer allows the serialization logic to be fully unit-tested without live API calls, making the test suite reliable and CI/CD-ready.

    \item \textbf{Bonus Intermediate Code Generation}: The TAC generator provides an early implementation of a compiler middle-end, laying the groundwork for future optimizations and machine code generation.
\end{enumerate}

\subsection{Summary of Deliverables as of Week 8}

\begin{tcolorbox}[colback=green!5!white, colframe=green!50!black, title={Final Deliverables Status}]
\begin{itemize}
    \item[\checkmark] \statusdone{COMPLETE} \textbf{EBNF Grammar} for Mini-Python (\texttt{grammar/mini\_python.ebnf})
    \item[\checkmark] \statusdone{COMPLETE} \textbf{Lexer} with indentation handling (\texttt{src/compiler/lexer.py})
    \item[\checkmark] \statusdone{COMPLETE} \textbf{Parser} with 16-node AST hierarchy (\texttt{src/compiler/parser.py}, \texttt{ast\_nodes.py})
    \item[\checkmark] \statusdone{COMPLETE} \textbf{Semantic Analyzer} with scope stack (\texttt{src/compiler/semantic\_analyzer.py})
    \item[\checkmark] \statusdone{COMPLETE} \textbf{Context Provider} (compiler $\to$ AI bridge) (\texttt{src/orchestrator/context\_provider.py})
    \item[\checkmark] \statusdone{COMPLETE} \textbf{Gemini Client} with system prompt + secure API (\texttt{src/orchestrator/gemini\_client.py})
    \item[\checkmark] \statusbonus{BONUS} \textbf{TAC Generator} (\texttt{src/compiler/tac\_generator.py})
    \item[\checkmark] \statusdone{COMPLETE} \textbf{128 Tests Passing} across Weeks 5--8
    \item[\checkmark] \statusdone{COMPLETE} \textbf{5 Demo Scripts} for each project phase
    \item[\checkmark] \statusdone{COMPLETE} \textbf{8 Weekly Reports} + SRS + Architecture documents (all in LaTeX)
\end{itemize}
\end{tcolorbox}


% ============================================================
\section{Conclusion}
% ============================================================

Project 83 demonstrates that a custom-built compiler front-end and a modern large language model can be integrated in a principled, structured way to create developer tools that are more intelligent, more explainable, and more secure than traditional compilers.

The key architectural insight --- that the compiler's internal state (\textbf{AST, Symbol Table, Diagnostics}) can be serialized into natural language and injected as context into an LLM prompt --- provides a powerful pattern applicable beyond this project: any structured reasoning system can be made AI-explainable by building a ``context bridge'' that translates its internal representations into human-readable text that an LLM can reason about.

With 128 tests passing and the complete compiler-to-AI pipeline operational, the project is well-positioned to complete the agentic tool-use, conversational interface, and security auditing phases in Weeks 9--16, delivering a full-featured Conversational Agentic AI Compiler Assistant.

\vspace{1cm}
\noindent\textbf{Prepared by:} Project 83 Team \\
\textbf{Document Date:} March 2026 \\
\textbf{Current Phase:} Week 8 Complete \\
\textbf{Next Milestone:} Week 9 --- Agentic Tool-Use (Function Calling)

\end{document}
